\paragraph{Theorem E.1.5 (Proved: integral homology of the fourfold).}
The fourfold of Theorem E.1.2 has torsion-free integral homology:
\[
 H_i(Z;\mathbb Z)=\mathbb Z^{b_i},\qquad
 (b_0,\ldots,b_8)=(1,0,2,1,2,1,2,0,1).
\]
In particular the finite-fibre torsion in E.1.4 does not survive
as torsion in the homology of the total space.

\begin{proof}
We compute the integral covering cokernels, rather than substituting
integral invariants in the rational calculation. Write
\(\mathcal V=R^1g_*\mathbb Z|_{B^\circ}\), and put
\(\mathcal J_q=j_*\bigwedge^q\mathcal V\).
The cohomology monodromies are the inverse transposes of
\(A_1\oplus C_1,A_2\oplus C_2,M_0\oplus C_0\).
All lattices below are in this marking.

First we describe a finite filling without making any equivariant
formality assumption. Since the first coordinate of each twist
\(v_j\) is \(\pm1\), the integral change of basis with first
column \(v_j/(v_j)_\gamma\) and remaining columns the last five
standard basis vectors changes its linear monodromy to
\(\operatorname{diag}(1,D_j)\) and its translation to
\(\pm e_0/m_j\). The reduced finite fibre is the mapping torus
of the actual linear automorphism \(D_j:T^5\to T^5\).
The covering \(T^6\to B_j\) is the \(m_j\)-fold cover of its
suspension circle. Set \(T_j=D_j^{-t}\). The two vertical matrices
on homology, in coordinates \((u,w,\delta,e_1,e_2)\), are
\[
D_1=\begin{pmatrix}
0&1&0&0&0\\-1&-1&0&0&0\\1&0&1&0&0\\
0&0&0&-1&-1\\0&0&0&1&0
\end{pmatrix},\qquad
D_2=\begin{pmatrix}
0&-1&0&0&0\\1&0&0&0&0\\0&1&1&0&0\\
0&0&0&0&1\\0&0&0&-1&0
\end{pmatrix}.
\]
The Wang sequences of the mapping torus and its cover give the
following exact description of the cokernel of covering pullback:
\[
 C_j^q\simeq
 \frac{\ker(\bigwedge^{q-1}T_j-I)}
 {\operatorname{im}\mathcal N_j^{q-1}},\qquad
 \mathcal N_j^{q-1}=\sum_{a=0}^{m_j-1}
                  (\bigwedge^{q-1}T_j)^a .
 \tag{E.I1}
\]
Here the target of pullback is the invariant lattice in
\(H^q(T^6;\mathbb Z)\). To verify (E.I1), its restriction to
the vertical torus surjects onto
\(\ker(\bigwedge^qT_j-I)\), and restriction of a mapping-torus
class has exactly the same image, by Wang exactness. Thus every
cokernel class has a representative containing the suspension
covector. On this kernel the covering map is the displayed norm:
cut the covered circle into its \(m_j\) consecutive intervals
and transport the vertical cochain through them. This proves both
surjectivity and the precise kernel of the horizontal description.
It does not choose a splitting of either Wang sequence.

For \(q=2\), the two quotients in (E.I1) are respectively zero
and \(\mathbb Z/2\). For \(q=3\), they are
\[
 C_1^3=\mathbb Z/3,\qquad
 C_2^3=(\mathbb Z/2)^2\oplus\mathbb Z/4.
 \tag{E.I2}
\]
The same Wang sequences show that the kernels of finite
specialization in degree two are \(\mathbb Z/3\) and
\(\mathbb Z/2\), and that the degree-three and degree-four
kernels vanish. The cusp specialization is integrally onto,
with the sole free degree-four kernel of E.1.3.

We give the integral global maps needed to use these local groups.
The global invariant two-form
\(q_X=u^*w^*+6\gamma\delta^*\) maps onto \(C_2^2\).
Indeed, if its pullback descended to the six-dimensional finite
fibre, restriction to the elliptic zero section would make it
descend to the original bielliptic surface \(S_2\). Its square
there would be \(\pm12/4=\pm3\), which is impossible in the
even intersection form of \(S_2\). That form is even because
\(H^2(S_2;\mathbb Z)\) is torsion free and its canonical class
is torsion, as checked in D.1.4. Thus the global-to-local map in
degree two is onto its only nonzero covering cokernel.

For degree three all three global invariant classes contain
\(\gamma\), and remain purely horizontal in both adapted
finite markings. They are
\[
 \gamma e_1^*e_2^*,\qquad
 \gamma(-u^*e_1^*+w^*e_2^*),\qquad
 \gamma u^*w^* .
 \tag{E.I3}
\]
Juxtaposition denotes exterior product. Consequently their images
in (E.I2) are exactly the vertical two-forms in (E.I3) modulo the
norm lattices; there is no unspecified lift in this computation.
In the integral invariant bases printed by the accompanying program,
the two norm matrices and this common image matrix are
\[
\begin{split}
N_3&=\begin{pmatrix}
0&0&0&0&0&0&0&0&0&3\\
0&0&-2&1&0&1&1&1&-1&0\\
0&0&-1&2&0&2&-1&0&-1&0\\
3&-1&0&0&2&0&0&0&0&0
\end{pmatrix},\\
N_4&=\begin{pmatrix}
0&0&0&0&0&0&0&0&0&4\\
0&0&-2&0&0&0&2&1&-1&0\\
0&0&0&2&0&2&0&-1&-1&0\\
4&-2&0&0&2&0&0&0&0&0
\end{pmatrix},\qquad
G=\begin{pmatrix}1&0&0\\0&1&0\\0&0&0\\0&0&1\end{pmatrix}.
\end{split}
\]
Every nonzero Smith factor of
\(\bigl(\operatorname{diag}(N_3,N_4)\mid(G,G)^t\bigr)\)
is one; there are eight. Thus the global classes surject onto
the direct sum of the two degree-three covering cokernels.

We also compute the integral parabolic groups, retaining torsion.
The degree-two global coboundary matrix, obtained by stacking
the two monodromy-minus-identity matrices, has a single nontrivial
Smith factor six. Its torsion generator restricts with orders
three, two and one at the three boundary circles. Hence no
nonzero torsion class is parabolic. Its rational parabolic rank
is zero, so \(H^1(B,\mathcal J_2)=0\). In degree three the
global coboundary image is primitive and the rational parabolic
rank is one; hence \(H^1(B,\mathcal J_3)=\mathbb Z\).
These assertions use the cocycle description as the kernel of
restriction from the free-group cohomology to the three
peripheral circle groups. The matrices and the restriction
orders are computed exactly in
\texttt{fourfolds\_integral\_audit.py}.

Passing from \(\mathcal J_q\) to \(R^qg_*\mathbb Z\)
adds the point-supported kernels just described and replaces
the finite stalks by the pullback sublattices. A point-supported
sheaf has no positive cohomology. The long exact sequences,
together with the two surjections onto covering cokernels,
therefore give
\[
 H^1(B,R^2g_*\mathbb Z)=0,\qquad
 H^1(B,R^3g_*\mathbb Z)=\mathbb Z.
 \tag{E.I4}
\]
They also identify degree-two base cohomology with the global
coinvariants. In the relevant degrees these are torsion-free,
with
\[
 H^2(B,R^1g_*\mathbb Z)=\mathbb Z,
 \qquad H^2(B,R^2g_*\mathbb Z)=\mathbb Z^2.
 \tag{E.I5}
\]
The first is the corresponding group from \(X\). The second
is the direct sum of that from \(\bigwedge^2\mathcal V_X\)
and the elliptic orientation line; the mixed coinvariants are
zero. These identifications, including primitivity, are also
verified by Smith forms, so no finite index is suppressed.

The integral differential onto the first group in (E.I5) is
surjective by naturality from the sphere fibration: its own
corresponding Leray quotient must vanish because \(H^3(X)=0\).
The differential onto the second group is surjective as well.
The \(X\) summand is killed by the differential transported
from \(H^4(X)=0\); the elliptic orientation summand is killed
by multiplying the primitive degree-one transgression of \(X\)
by the permanent class of the elliptic zero-section divisor.
That divisor has degree one on each elliptic fibre. Thus these
are two primitive independent target generators, not merely two
independent rational images.

Now \(H^0(B,R^3g_*\mathbb Z)\) and
\(H^0(B,R^4g_*\mathbb Z)\) are free: in the first degree
specialization has no kernel, and in the second its only
kernel is the free cusp class. There is no later Leray
differential, since the base has cohomological dimension two.
Equation (E.I4) and the first surjectivity in (E.I5) show that
\(H^3(Z;\mathbb Z)\) is a free group, of rank one by E.1.4.
Equation (E.I4) and the second surjectivity show that
\(H^4(Z;\mathbb Z)\) is an extension of two free groups of
rank one, and hence is \(\mathbb Z^2\).
Simple connectivity makes \(H^2(Z;\mathbb Z)\) torsion free.
The universal coefficient theorem now proves successively that
\(H_2(Z;\mathbb Z)=\mathbb Z^2\) and
\(H_3(Z;\mathbb Z)=\mathbb Z\). Poincare duality determines
the other degrees and gives exactly the displayed groups.
\end{proof}
